\documentclass[a4paper]{article}
% Espanol (Mexico) con XeLaTeX: fontspec para fuentes Unicode, polyglossia
% para la division silabica y los nombres del documento en la variante mexicana.
\usepackage{fontspec}
\usepackage{polyglossia}
\setdefaultlanguage[variant=mexican]{spanish}
% polyglossia no traduce los nombres de listings.
\AtBeginDocument{\providecommand{\lstlistingname}{}\renewcommand{\lstlistingname}{Listado}%
  \providecommand{\lstlistlistingname}{}\renewcommand{\lstlistlistingname}{Índice de listados}}
\usepackage{amsmath,amssymb}
\usepackage{graphicx}
\usepackage[margin=2.35cm]{geometry}
\usepackage[most]{tcolorbox}
\usepackage{etoolbox}
\usepackage{listings}
\usepackage{xcolor}
\usepackage{hyperref}
\usepackage{booktabs,longtable,tabularx,array}
\usepackage{subcaption}
\usepackage{float}
\usepackage{enumitem}
\usepackage{microtype}
\usepackage{fancyhdr}
\usepackage{needspace}

\definecolor{kbBack}{HTML}{F2F6FA}
\definecolor{kbFrame}{HTML}{9DB4CC}
\definecolor{kbTitle}{HTML}{52708F}
\definecolor{ibBack}{HTML}{FBF5E7}
\definecolor{ibFrame}{HTML}{C9A961}
\definecolor{ibTitle}{HTML}{7D5F1E}
\definecolor{wbBack}{HTML}{FCF2F0}
\definecolor{wbFrame}{HTML}{D6A096}
\definecolor{wbTitle}{HTML}{9E4F43}
\definecolor{dbBack}{HTML}{F4F7F3}
\definecolor{dbFrame}{HTML}{AEC2AC}
\definecolor{dbTitle}{HTML}{5F7F64}

\tcbset{softbox/.style={
  enhanced,breakable,boxrule=0.8pt,arc=2pt,left=3mm,right=3mm,
  top=2mm,bottom=2mm,coltitle=white,fonttitle=\bfseries,
  toptitle=0.6mm,bottomtitle=0.6mm
}}
\newtcolorbox{knowledgebox}[1]{softbox,colback=kbBack,colframe=kbFrame,colbacktitle=kbTitle,title={#1}}
\newtcolorbox{importantbox}[1]{softbox,colback=ibBack,colframe=ibFrame,colbacktitle=ibTitle,title={#1}}
\newtcolorbox{warningbox}[1]{softbox,colback=wbBack,colframe=wbFrame,colbacktitle=wbTitle,title={#1}}
\newtcolorbox{dialoguebox}[1]{softbox,colback=dbBack,colframe=dbFrame,colbacktitle=dbTitle,title={#1}}

\lstset{
  language=Python,basicstyle=\ttfamily\small,
  keywordstyle=\color{kbTitle}\bfseries,stringstyle=\color{wbTitle},
  commentstyle=\color{dbTitle},breaklines=true,frame=single,numbers=left,
  numberstyle=\tiny\color{gray},captionpos=b,columns=fullflexible,
  keepspaces=true,showstringspaces=false,extendedchars=true
}
\BeforeBeginEnvironment{lstlisting}{\Needspace{12\baselineskip}}

\hypersetup{colorlinks=true,linkcolor=kbTitle,urlcolor=kbTitle,citecolor=wbTitle}
\setlist{itemsep=0.25em,topsep=0.4em}
\setlength{\parindent}{2em}
\setlength{\parskip}{0.25em}
\newcolumntype{L}[1]{>{\raggedright\arraybackslash}p{#1}}

\providecommand{\notetitle}{Notas del curso: [tema del curso]}
\providecommand{\noteauthors}{Elaboradas a partir de los materiales públicos de Stanford CS153}
\providecommand{\notedate}{Primavera de 2026}
\providecommand{\videochannel}{Stanford Online}
\providecommand{\videopublishdate}{2026}
\providecommand{\videoduration}{[duración del video]}
\providecommand{\videourl}{}
\providecommand{\videocoverpath}{cover.jpg}
\providecommand{\coursesite}{https://cs153.stanford.edu/}
\providecommand{\repourl}{https://github.com/hqhq1025/ai-course-notes}
\providecommand{\skillversion}{wdkns/wdkns-skills@39f1a04}

\newcommand{\figsource}[1]{\par\vspace{-0.3em}{\footnotesize\color{gray}\emph{Fuente: #1}}\par}
\newcommand{\teachervoice}[1]{\begin{knowledgebox}{Nota de clase}#1\end{knowledgebox}}
\newcommand{\term}[1]{\textbf{#1}}

\pagestyle{fancy}
\fancyhf{}
\fancyfoot[C]{\thepage}
\fancyfoot[R]{\footnotesize\color{gray}\href{\repourl}{AI Course Notes}}
\renewcommand{\headrulewidth}{0pt}
\renewcommand{\footrulewidth}{0pt}

\newcommand{\makecscover}{
\begin{titlepage}
\centering
\vspace*{0.7cm}
{\Large Stanford CS153 · Frontier Systems\par}
\vspace{0.8cm}
{\huge\bfseries \notetitle\par}
\vspace{0.6cm}
{\large \noteauthors\par}
\vspace{0.25cm}
{\large \notedate\par}
\vspace{0.9cm}
\ifdefempty{\videocoverpath}{}{
  \includegraphics[width=0.86\textwidth,height=0.43\textheight,keepaspectratio]{\videocoverpath}\par
}
\vfill
\begin{tcolorbox}[width=0.92\textwidth,colback=black!2!white,colframe=black!45,boxrule=0.7pt,arc=2pt]
\textbf{Autor o canal del video}: \videochannel\par
\textbf{Fecha de publicación}: \videopublishdate\par
\textbf{Duración del video}: \videoduration\par
\textbf{Sitio del curso}: \href{\coursesite}{\nolinkurl{\coursesite}}\par
\ifdefempty{\videourl}{}{\textbf{Enlace del video}: \href{\videourl}{\nolinkurl{\videourl}}\par}
\textbf{Normas de elaboración}: \skillversion\ + las normas source-first / teacher-voice / visual-QA de este repositorio
\end{tcolorbox}
\end{titlepage}
\tableofcontents
\newpage
}


\renewcommand{\notetitle}{Lecture 11: Security Leadership — VDP, respuesta a incidentes, gobernanza de la divulgación y responsabilidad de los altos directivos}
\renewcommand{\noteauthors}{Elaborado a partir de la entrevista a Joe Sullivan, los subtítulos con marcas de tiempo y fuentes oficiales del gobierno y los tribunales de Estados Unidos}
\renewcommand{\notedate}{Curso de invierno de 2025 · Versión reescrita en 2026 con enfoque source-first}
\renewcommand{\videochannel}{Publicación de cursos anteriores de Stanford CS153}
\renewcommand{\videopublishdate}{2025-01-23}
\renewcommand{\videoduration}{39:17}
\renewcommand{\videourl}{https://www.youtube.com/watch?v=u4E6qdWIkYI}
\renewcommand{\videocoverpath}{cover.jpg}

\newcommand{\lecturefigure}[3]{%
\begin{figure}[H]
  \centering
  \includegraphics[width=0.96\textwidth]{images/#1}
  \caption{#2}
  \figsource{#3}
\end{figure}
}

\begin{document}
\makecscover

\section{Auditoría de fuentes: convertir un caso personal en una clase de ingeniería organizacional}

La trayectoria profesional de Joe Sullivan abarca la fiscalía federal de delitos informáticos, la security leadership en empresas de tecnología, el ecosistema de divulgación de vulnerabilidades y proyectos de interés público. Por ello, la clase combina experiencia profesional, el relato de un caso personal, juicios de política pública y recomendaciones de liderazgo. Unos buenos apuntes no pueden mezclar estos contenidos como si fueran un solo tipo de «hecho»: el recuerdo que el ponente tiene de los acontecimientos es una declaración hecha en clase, las afirmaciones de la fiscalía son acusaciones, el jurado y los tribunales se ocupan de la responsabilidad jurídica, y este texto se encarga de extraer mecanismos de ingeniería y de gobernanza reutilizables.

\begin{warningbox}{Límites jurídicos y temporales}
Este texto son apuntes de security engineering y governance; no constituye asesoría jurídica. La clase se grabó en enero de 2025, cuando la apelación de Sullivan aún no se resolvía; el Tribunal de Apelaciones del Noveno Circuito confirmó la condena el 13 de marzo de 2025, el 12 de noviembre de 2025 publicó una opinión modificada y negó la revisión en pleno, y la Suprema Corte de Estados Unidos se negó a admitir la certiorari petition el 29 de junio de 2026. Las expectativas sobre la apelación expresadas en clase solo pueden tomarse como una opinión histórica, no como el estado procesal actual.
\end{warningbox}

\begin{importantbox}{Los tres ciclos de gobernanza de esta clase}
\begin{enumerate}
  \item \textbf{External-report loop}: recibir el reporte externo $\rightarrow$ verificarlo y clasificar su gravedad $\rightarrow$ corregir $\rightarrow$ divulgación coordinada o recompensa;
  \item \textbf{Incident loop}: detectar $\rightarrow$ contener el impacto $\rightarrow$ preservar la evidencia $\rightarrow$ investigar, recuperar y hacer la revisión posterior;
  \item \textbf{Disclosure loop}: hechos técnicos $\rightarrow$ análisis jurídico y de materiality $\rightarrow$ decisión de la alta dirección o del consejo de administración $\rightarrow$ comunicación a reguladores, clientes o al público.
\end{enumerate}
Los tres ciclos pueden detonarse a partir de una misma pista, pero no pueden sustituirse por una misma persona, un mismo ticket ni un mismo mecanismo de incentivos.
\end{importantbox}

\subsection{Un security program no es un «CSO héroe»}

Con frecuencia se evalúa a los líderes de seguridad por su capacidad individual: si son capaces de detectar un ataque, de convencer al CEO o de decidir con rapidez en plena crisis. Sin embargo, lo que en realidad necesita una organización a escala es un \term{security program} (programa de seguridad): capacidades del personal bien definidas, procesos repetibles, plataformas técnicas y una asignación clara de facultades y responsabilidades de gobernanza, de modo que la conducta correcta no dependa de la improvisación de un responsable. De la transición de Sullivan de fiscal a empresas de tecnología como eBay, Facebook y Uber, el tema común más valioso no son los nombres de las empresas, sino que el equipo de seguridad debe evolucionar al mismo ritmo que la complejidad de la organización.

\lecturefigure{01-security-program.png}{Un security program sostenible depende a la vez de people, process, platform y governance.}{Subtítulos locales 00:00--03:20; reelaboración conceptual.}

\begin{knowledgebox}{Lectura de la figura: ninguna de las cuatro capas puede faltar}
Primero observe people, en la esquina superior izquierda: el talento, las guardias, la comunicación y la capacidad de decisión determinan si la organización logra comprender un incidente. Process convierte la intake, el incident command, la disclosure y la revisión posterior en rutas repetibles. Platform aporta identity, telemetry, case management y evidence storage. Governance responde quién tiene authority, quién se somete al challenge y quién firma la divulgación externa. Si solo se construyen herramientas, el resultado es «muchas alertas y pocas decisiones»; si solo se depende de expertos, el resultado es «el experto se va y el sistema pierde la memoria».
\end{knowledgebox}

\teachervoice{La historia profesional de Sullivan reúne el derecho, la ingeniería y la organización en una misma perspectiva. Nota de clase: el valor escaso de un security leader no consiste en conocer cada vulnerabilidad mejor que todos los ingenieros, sino en traducir los hechos técnicos a riesgos jurídicos, de negocio y públicos, y al mismo tiempo asegurar que cada juicio lo emita quien tiene la facultad de responder por él.}

\begin{knowledgebox}{Términos clave: funciones y responsabilidades}
El \term{CSO/CISO} es responsable de las capacidades de seguridad, de la información sobre riesgos y del escalamiento; el \term{general counsel} (consejero jurídico general, GC) es responsable del análisis jurídico de la organización y de la gestión de los abogados; el \term{CEO} responde por las prioridades de la empresa y por las decisiones de negocio importantes; el \term{board/disclosure committee} ejerce la supervisión y atiende las divulgaciones relevantes. Las funciones pueden fusionarse según el tamaño de la empresa, pero debe distinguirse de forma explícita «quién aporta los hechos, quién emite la opinión jurídica y quién asume la decisión de negocio».
\end{knowledgebox}

\subsection{Resumen de la sección}

Esta clase no gira en torno a «quién tiene más razón», sino en torno a un operating system: el reporte externo, la respuesta a incidentes y la gobernanza de la divulgación requieren, cada uno, owner, evidence, decision right y registro. El caso personal solo sirve para comprobar por qué estos mecanismos tienen que existir.
\section{Tecnología, gobierno y reglas: por qué los ciclos de retroalimentación siempre llegan tarde}

Las empresas tecnológicas convierten en sistemas de software funciones sociales como los pagos, las comunicaciones, la identidad, el transporte, la salud y la IA, y con ello acumulan un poder que antes ejercían instituciones públicas o infraestructuras reguladas. El gobierno debe responder a los daños, pero a menudo va rezagado en talento técnico, visibilidad de los datos y velocidad para iterar sus reglas. Si las empresas ven la regulación como una fricción puramente externa, pierden la oportunidad de dar forma a reglas que puedan cumplirse; si el gobierno solo aplica la ley después de los incidentes, también genera ambigüedad y un efecto inhibidor.

\subsection{Public-private feedback loop}

El ciclo de política pública ideal no es «las empresas innovan primero y el gobierno sanciona al final», sino uno en el que las señales tecnológicas, los daños reales, el diseño de reglas, la capacidad de implementación y la retroalimentación de la industria se corrigen mutuamente de forma continua. El problema de ingeniería central aquí es la observabilidad: el gobierno no ve los riesgos internos de las plataformas y las empresas no conocen las restricciones de la política pública, de modo que cada parte diseña controles con información incompleta.

\lecturefigure{02-public-private-loop.png}{La tecnología y el gobierno necesitan retroalimentación continua, no una rendición de cuentas unidireccional después de los incidentes.}{Subtítulos locales 03:20--06:20; redibujo conceptual.}

\begin{knowledgebox}{Lectura de la figura: la retroalimentación debe volver a la capa de implementación}
Technology produce productos, datos y riesgos nuevos; los observed harms generan señales de los usuarios, de la aplicación de la ley y del mercado; policy plasma los objetivos en leyes, estándares o expectativas regulatorias; implementation los lleva a controls, reporting y feedback. Si el ciclo se detiene en el texto de la política, la organización solo obtiene obligaciones abstractas; si se detiene en la autorregulación empresarial, el público afectado no cuenta con una reparación confiable. Una buena retroalimentación convierte las reglas en controles que pueden probarse y, a la vez, devuelve al diseño de la política los fallos de esos controles.
\end{knowledgebox}

\begin{warningbox}{Tanto «no regular» como «cualquier regulación es buena» son posturas demasiado simples}
Las reglas pueden quedar obsoletas, superponerse o premiar el cumplimiento meramente formal, pero la falta de reglas claras también traslada los costos a los usuarios, a los investigadores y a los equipos de seguridad de primera línea. Lo que el curso pide de verdad es participar en el diseño: definir con claridad a quién se quiere proteger, qué evidencia hay que observar, qué margen de experimentación se permite, cuáles son los umbrales de escalamiento y cuáles los mecanismos de apelación, en lugar de limitarse a expresar apoyo u oposición.
\end{warningbox}

\teachervoice{El ponente enfatiza que, en sus inicios, la industria tecnológica solía repetir como un mantra que «la regulation asfixia la innovación». El giro más maduro que se plantea en clase es que la falta de reglas también genera fricción: las empresas no conocen los límites, los usuarios no tienen confianza y el gobierno solo puede definir los estándares a posteriori, a partir de incidentes y de la aplicación de la ley.}

\subsection{Por qué canales llegan las reglas a las organizaciones}

Las obligaciones de seguridad rara vez provienen de una sola disposición legal. El Congress o las legislaturas estatales pueden crear obligaciones, las agencias administrativas pueden hacer rulemaking o ejercer supervisión, organismos de normalización como el NIST ofrecen marcos operativos, y la aplicación de la ley y los tribunales aplican las reglas generales a hechos concretos. Las empresas multinacionales enfrentan además distintas jurisdicciones, reglas sectoriales, contratos y compromisos con clientes. El sistema de gobernanza debe mantener un obligation inventory, en lugar de que cada incident obligue a buscar sobre la marcha «qué podría aplicar».

\lecturefigure{11-regulation-paths.png}{Las reglas de ciberseguridad se forman por varias vías: legislación, estándares, supervisión, aplicación de la ley y tribunales.}{Subtítulos locales 06:00--09:10; documentos gubernamentales; redibujo conceptual.}

\begin{knowledgebox}{Lectura de la figura: distinguir el origen de la regla y el tipo de evidencia}
Legislation/rules define obligaciones generales; los standards traducen los objetivos a controls y practices; la supervision observa el cumplimiento mediante inspecciones, informes e interacción continua; enforcement/courts emiten juicios sobre expedientes concretos. Los equipos de ingeniería no pueden equiparar un estándar con la ley, ni tomar un solo caso de aplicación de la ley como una norma completa. Lo correcto es que el counsel establezca el juicio de aplicabilidad y que el equipo de seguridad mapee las obligaciones a activos, procesos, registros y evidencia de simulacros.
\end{knowledgebox}

\begin{importantbox}{Obligation-to-control mapping}
Para cada obligación se mantienen cinco columnas: \textbf{source} (origen), \textbf{scope} (qué entidades, datos o sistemas), \textbf{trigger} (cuándo entra en vigor), \textbf{control/evidence} (cómo se demuestra el cumplimiento) y \textbf{owner/reviewer} (quién es responsable y quién lo cuestiona). Así, el incident commander puede ubicar rápidamente los requisitos pertinentes, en lugar de cargarle la investigación jurídica al on-call engineer en un momento de alta presión.
\end{importantbox}

\subsection{Resumen de la sección}

La tensión entre la tecnología y el gobierno es, en esencia, un problema de retroalimentación, capacidad y evidencia. Las organizaciones deben mapear de forma proactiva las leyes, los estándares y la experiencia en aplicación de la ley a controles ejecutables y, al mismo tiempo, hacer que el personal técnico participe en la discusión de políticas, para que las reglas reflejen cómo fallan realmente los sistemas.
\section{Security Team Maturity: del equipo de emergencias a la gobernanza integrada}

La madurez de un equipo no es una curva de headcount. Un equipo de 5 personas puede tener una authority clara y una automatización de alta calidad; un equipo de 100 personas puede limitarse a recibir alertas y reenviar tickets. Una security organization madura integra de forma gradual el threat signal, el incident workflow, la risk acceptance y la product decision en los sistemas del negocio, de modo que la seguridad deja de ser una aprobación provisional antes de cada lanzamiento.

\subsection{Cuatro etapas y los atajos equivocados}

La etapa reactive depende de unos cuantos expertos que atienden emergencias; la etapa program establece intake, on-call, policy y repeatable controls; la etapa platform convierte identity, telemetry, automation y guardrails en capacidades compartidas; la etapa embedded hace que product, finance, legal, HR y communications asuman su propia responsabilidad de seguridad en las decisiones cotidianas. Las etapas no avanzan por sí solas y, sobre todo, la compra de productos no puede usarse para aparentar que se completó un cambio organizacional.

\lecturefigure{13-security-maturity.png}{El equipo de seguridad pasa de la reactive response a la embedded governance.}{Subtítulos locales 00:00--03:20; 31:20--33:40; concepto redibujado.}

\begin{knowledgebox}{Lectura de la figura: el resultado de la madurez no son más controles, sino ciclos de retroalimentación más cortos}
Reactive solo puede responder después de un incidente; program establece roles y un ritmo de operación; platform hace que los controles sean reutilizables y observables; embedded devuelve al negocio la ownership del producto, la risk acceptance y la mejora continua. Para evaluar la madurez, primero se revisa si los handoff están bien definidos, si las decisiones dejan registro, si las excepciones caducan y si el postmortem modifica el sistema, no la cantidad de policy ni el monto de las compras.
\end{knowledgebox}

\begin{warningbox}{«Security owns risk» es una formulación peligrosa}
El equipo de seguridad suele encargarse de identificar, analizar, recomendar y monitorear, pero los responsables del negocio son dueños de los beneficios del producto y del riesgo residual, el área legal se encarga del juicio jurídico, y los directivos y el consejo de administración asumen las decisiones empresariales de mayor peso. Entregar todo el riesgo al CSO hace que los demás responsables estén ausentes en tiempos normales y que, tras un incidente, se busque a un único culpable.
\end{warningbox}

\begin{importantbox}{La interfaz mínima de un equipo maduro}
Como mínimo, hay que definir: intake unificado, severity taxonomy, incident commander, legal/comms escalation, evidence repository, materiality workstream, exception register, executive exercise y post-incident owner. Cada interfaz debe tener un tiempo de respuesta, un rol suplente y salidas que puedan auditarse.
\end{importantbox}

\subsection{Resumen de la sección}

El objetivo del crecimiento del equipo es codificar la experiencia individual en interfaces organizacionales. Cuanto más cerca se está de la etapa embedded, menos aparecen las rupturas del tipo «el equipo de seguridad lo sabe, pero los directivos no», «el área legal emitió su juicio, pero sin los hechos técnicos» o «el incidente terminó, pero nadie modificó el sistema».
\section{Vulnerability Disclosure Policy: convertir la investigación externa en insumo para las correcciones}

Los investigadores externos pueden descubrir problemas antes que el equipo interno, pero si no existe un canal de reporte claro, no saben qué sistemas pueden probar, qué conductas están permitidas ni cuándo recibirán respuesta. La \term{Vulnerability Disclosure Policy} (VDP, política de divulgación de vulnerabilidades) es un conjunto público de reglas y un mecanismo de recepción: define el scope, los límites de la autorización, la forma de envío, los compromisos de la organización y la comunicación coordinada. Lo primero que resuelve es «cómo reportar de forma segura»; no garantiza el pago de una recompensa.

\subsection{Del report a la coordinated disclosure}

La \term{coordinated vulnerability disclosure} (divulgación coordinada de vulnerabilidades, CVD) es el proceso en el que los investigadores, la organización afectada y los posibles proveedores coordinan, entre la corrección y la divulgación pública, los tiempos, el alcance del impacto y el contenido de la información. El valor de ingeniería de una VDP está en convertir correos ambiguos en un workflow rastreable: acknowledge, triage, validate, scope, remediate, verify, communicate. Aquí el service-level objective no sirve para tener un dashboard vistoso, sino para reducir la probabilidad de que los investigadores repitan pruebas, escalen el asunto públicamente o pierdan la confianza.

\lecturefigure{03-vdp-intake.png}{La VDP convierte los reportes de investigación externa en un proceso de validación, corrección y divulgación coordinada.}{CISA VDP template; subtítulos locales 12:30--15:20; rediseño conceptual.}

\begin{knowledgebox}{Lectura de la figura: cada handoff debe devolver un estado}
El reporter envía el activo, el impacto y la evidencia de reproducción; el triage confirma el scope, la severity y si es duplicate; la remediation asigna owner, corrección, verificación y timeline; el close/disclose actualiza el estado ante el investigador y, cuando hace falta, coordina con la cadena de suministro o con el público. Lo que falla con más facilidad es el silencio intermedio: la organización ya recibió el reporte pero no lo confirma, ya corrigió pero no verifica, necesita una prórroga pero no da razones. Transparencia no significa publicar toda la información interna, sino establecer compromisos predecibles.
\end{knowledgebox}

\begin{importantbox}{Campos mínimos del intake de una VDP}
\begin{itemize}
  \item asset, endpoint o product version afectados;
  \item security impact observado y evidencia mínima de reproducción;
  \item fecha y hora de la prueba, researcher contact y si podrían estar involucrados terceros;
  \item case ID generado automáticamente, acknowledgement y ventana de actualización de estado;
  \item severity, owner, containment, fix, verification y disclosure decision.
\end{itemize}
No pidas a los investigadores datos personales ajenos a la validación, ni dejes que el bounty payment sea la única prioridad para programar las correcciones.
\end{importantbox}

\teachervoice{Sullivan describe la responsible disclosure como una innovación de infraestructura: les dice a los investigadores de buena fe «si reportas dentro de estos límites, queremos saberlo y lo atenderemos por los canales establecidos». Nota de clase: se trata de diseñar una relación, no solo de publicar una dirección en security.txt.}

\begin{warningbox}{La VDP no puede degradar automáticamente un incidente a vulnerabilidad común}
Un reporte externo puede indicar al mismo tiempo unauthorized access, credential compromise, obtención de datos o una amenaza persistente. En ese caso, el caso de VDP sigue a cargo de la researcher communication, pero el incidente debe pasar a un incident command independiente, con evidence preservation y disclosure analysis. Dejar todos los hechos en la bounty platform oculta la gravedad.
\end{warningbox}

\subsection{Resumen de la sección}

La VDP es un canal de reporte de seguridad y un contrato de confianza. Mediante el scope, los compromisos de respuesta y la corrección coordinada, reduce la incertidumbre de ambas partes; pero en cuanto los hechos muestran un compromise real, debe activar una gobernanza de incidentes más sólida.
\section{Bug Bounty y Safe Harbor: la recompensa no sustituye los límites de la autorización}

Un VDP y un bug bounty suelen aparecer juntos, pero resuelven problemas distintos. Un \term{bug bounty program} (programa de recompensas por vulnerabilidades) es un programa opcional que, bajo reglas establecidas, ofrece dinero, puntos o reconocimiento a los reportes válidos que cumplen las condiciones; un VDP debe existir aunque no haya recompensa. El bounty ayuda a competir por la atención de los investigadores, a fijar severity rewards y a operar el triage, pero un pago posterior no puede reescribir una conducta dañina que ya ocurrió.

\subsection{El límite entre VDP y bounty}

Una organización puede tener un VDP sin bounty, y también puede administrar el bounty mediante una plataforma de terceros. Ambos comparten el alcance de activos, el canal de reporte y el mecanismo de coordinación, pero el bounty requiere además reward table, duplicate rule, eligibility, tax/payment y appeal. Los responsables de seguridad deben evitar que el equipo pase por alto data access, extortion signal u otros incident indicators porque «es un reporte de bounty».

\lecturefigure{04-vdp-vs-bounty.png}{El VDP ofrece un canal seguro de reporte; el bug bounty agrega sobre él un mecanismo opcional de recompensa.}{CISA/DOJ VDP guidance; subtítulos locales 12:30--15:20; redibujado conceptual.}

\begin{knowledgebox}{Lectura de la figura: primero la autorización, después la recompensa}
El núcleo del VDP es safe channel, scope, response promise y coordinated disclosure; la capa adicional del bounty es optional payment, severity table, eligibility y duplicate. Ninguno de los dos cubre la extortion, el sabotaje, el acceso a datos más allá de lo necesario ni la negativa a detener las pruebas. Al evaluar un reporte, primero se determina si la conducta se ajustó a los límites publicados de antemano y a la good faith, y después se valoran la vulnerabilidad y la recompensa; el proceso de pago no debe decidir la clasificación legal ni la del incidente.
\end{knowledgebox}

\begin{warningbox}{El payment no es una post-hoc authorization}
Que la organización pague una recompensa, firme un NDA o mantenga la comunicación puede servir para controlar el riesgo y facilitar la corrección; estas acciones no prueban por sí mismas que la conducta previa estuviera autorizada, ni eliminan la necesidad de conservar evidencia, notificar a los afectados o consultar a counsel. Cualquier atajo del tipo «pagar para convertir el incident en bounty» crea un punto ciego de gobernanza.
\end{warningbox}

\subsection{Safe harbor, scope y good faith}

En el contexto de un VDP, el \term{safe harbor} (puerto seguro) es el compromiso de la organización con la investigación de buena fe que cumple la política; por ejemplo, no ejercer ciertas acciones civiles, respaldar al investigador para que demuestre su buena fe o, dentro del ámbito aplicable, no recomendar acciones penales. No es una inmunidad legal general ni puede obligar a todos los terceros. El \term{scope} (alcance) especifica los dominios, aplicaciones, API, versiones y activos excluidos que pueden probarse; la \term{good faith} (buena fe) suele exigir solo el acceso mínimo necesario para la verificación, evitar afectaciones a la privacidad y al servicio, y detenerse y reportar a tiempo.

\lecturefigure{05-safe-harbor.png}{El safe harbor solo es eficaz si se combina con un scope claro, requisitos de buena fe, conductas prohibidas y una ruta de escalamiento.}{CISA VDP template; DOJ VDP framework; redibujado conceptual.}

\begin{knowledgebox}{Lectura de la figura: la política debe proteger a la vez a los investigadores y a los usuarios}
In-scope enumera los activos y métodos que se permite probar; good faith exige minimizar el acceso, reportar y proteger los datos sensibles; prohibited fija límites como disruption, persistence, data use o extortion; escalation ofrece un safety contact, un legal contact y una vía para las disputas. Si solo hay frases de aliento sin un scope concreto, el investigador asume riesgos impredecibles; si solo hay una lista de prohibiciones sin compromiso de respuesta, no se genera confianza.
\end{knowledgebox}

\begin{importantbox}{Lista de verificación para diseñar la política}
El inventario de activos debe poder mantenerse; los servicios de terceros deben indicar los límites de propiedad; las restricciones de velocidad de prueba y de automatización deben explicar su motivo; el manejo de datos sensibles debe establecer cuándo detenerse, la retención mínima y la eliminación; las disputas deben contar con un escalation humano. Las actualizaciones de la política también deben conservar la versión y la fecha de entrada en vigor, para que la organización no evalúe conductas anteriores con reglas nuevas.
\end{importantbox}

\subsection{Resumen de la sección}

El VDP se ocupa del canal, el bounty de los incentivos y el safe harbor del compromiso de la organización con quien respeta los límites. Ninguno de los tres sustituye la incident classification, y ninguno permite tratar el pago, el NDA o la presunción de buena fe como hechos técnicos.
\section{Incident Command: cuándo un reporte externo escala a incidente corporativo}

Cuando el reporte muestra que se usó una credential, que se accedió a sistemas sensibles, que los datos pudieron salir del perímetro de control o que el ataque sigue en curso, lo que enfrenta la organización ya no es solo vulnerability remediation. El \term{incident commander} (comandante del incidente, IC) es el único responsable operativo que coordina objetivos, prioridades, ritmo, roles y escalamiento. El IC no tiene que hacer personalmente el análisis forense ni emitir conclusiones jurídicas, pero debe asegurar que cada workstream tenga owner, que los conflictos se escalen y que las decisiones queden registradas.

\subsection{Estructura de mando interfuncional}

La respuesta a incidentes incluye al menos cuatro workstreams: command, technical, legal/comms y business. El equipo technical confirma los hechos, controla el impacto, corrige y restablece; legal/comms identifica obligaciones, privilege y audiencias externas; business resuelve las decisiones entre usuarios, operación, finanzas y producto; command mantiene la timeline, los objectives, la cadence y la dependency comunes. Si el security on-call asume todos los roles a la vez, bajo presión terminará registrando supuestos como si fueran hechos.

\lecturefigure{06-incident-command.png}{Un incidente de seguridad grave requiere una estructura interfuncional de command, technical, legal/comms y business.}{NIST SP 800-61 Rev. 3; subtítulos locales 15:10--25:40; reconstrucción conceptual.}

\begin{knowledgebox}{Lectura de la figura: el IC sincroniza, no monopoliza el juicio especializado}
Command mantiene la severity, los owners y el decision log; technical se ocupa de contain, forensics y recovery; legal/comms analiza privilege, regulators, customers y messaging; business decide sobre customer impact, continuity y recursos. Las flechas son bidireccionales: un hallazgo técnico puede cambiar la materiality, un deadline legal puede cambiar las prioridades de la investigación y las acciones de recuperación del negocio pueden destruir evidencia; por eso ningún flujo de trabajo puede operar de forma aislada.
\end{knowledgebox}

\begin{knowledgebox}{Términos clave: severity, containment y recovery}
\term{Severity} es la clasificación operativa que hace la organización del impacto y la urgencia actuales; no equivale a la materiality jurídica final. \term{Containment} es la acción de corto plazo que limita el ataque o su impacto sobre los datos, por ejemplo aislar una credential o un servicio. \term{Eradication} elimina la causa raíz y la persistence; \term{recovery} restablece un servicio confiable y vigila que el problema no se repita. Restablecer primero el negocio no significa que la investigación haya terminado, y preservar primero la evidencia tampoco obliga a aplazar todas las acciones que protegen a los usuarios.
\end{knowledgebox}

\teachervoice{De la exposición en clase, lo que más vale la pena conservar no son los pasos concretos de la intrusión, sino el conflicto de clasificación: una pista puede llegar al principio por el disclosure channel y revelarse poco a poco como un incidente más grave. La lección práctica es definir de antemano las condiciones de escalamiento, para que el personal no tenga que decidir por su cuenta con base en «¿esto parece o no una bounty?».}

\begin{importantbox}{VDP-to-incident escalation triggers}
Se inicia de inmediato una incident review cuando se detecta uso de cuentas reales o de credentials, acceso a datos no minimizado, salida de datos sensibles fuera del perímetro, acceso persistente, afectación a terceros, una exigencia de pago del investigador sin relación con detener el daño, una discrepancia evidente entre los hechos y el reporte, o cualquier trigger regulatorio o contractual. Que se active un trigger no equivale a acusar de un delito al investigador: significa elevar el nivel de investigación y de gobernanza de la organización.
\end{importantbox}

\subsection{Resumen de la sección}

El mando del incidente organiza varios juicios especializados en un ritmo operativo común. El valor del IC está en que los hechos, las acciones de protección, las obligaciones legales y las decisiones de negocio sean visibles entre sí, y no en que el CSO haga a la vez de ingeniero, abogado, vocero y CEO.
\section{Evidence Preservation: que los hechos puedan verificarse}

En un incidente grave, la memoria del equipo se altera por la presión, las conversaciones paralelas y los resultados posteriores. La preservación de evidencia no busca «preparar una versión para un litigio», sino que la investigación técnica, la protección de los usuarios, los seguros, la respuesta a reguladores y la revisión posterior partan del mismo conjunto de insumos verificables. La preservación debe comenzar lo antes posible y, al mismo tiempo, no debe obstaculizar la containment.

\subsection{Chain of custody, legal hold y timeline}

La \term{chain of custody} (cadena de custodia) registra quién recopiló, transfirió, consultó y almacenó la evidencia, cuándo y de qué manera, de modo que un revisor posterior pueda juzgar su integridad. Un hash ayuda a detectar cambios en un archivo, pero el hash por sí mismo no demuestra que el método de adquisición fuera correcto. Un \term{legal hold} (aviso de conservación legal) es el procedimiento por el cual el counsel ordena suspender la eliminación habitual y conservar la información que pueda ser pertinente; los ingenieros no deben declarar por su cuenta el privilege ni decidir el alcance de la conservación.

\lecturefigure{07-evidence-chain.png}{La calidad de la evidencia depende de una cadena continua de preserve, integrity, custody y timeline.}{Guía de respuesta a incidentes del NIST; subtítulos locales 21:00--25:40; redibujo conceptual.}

\begin{knowledgebox}{Lectura de la figura: la integridad proviene del proceso, no solo de las herramientas}
La etapa de preserve determina los logs, systems, snapshots y communications; integrity registra el hash, el timestamp y el acquisition method; custody restringe y audita el acceso, la transfer y el purpose; timeline ordena en el tiempo los hechos, las hipótesis, las decisiones y las correcciones posteriores. No basta con tomar una captura de pantalla de una conversación, ni debe permitirse que todos modifiquen directamente el mismo directorio de evidencia. Cada exportación debe tener un owner, un sistema de origen, una zona horaria y una política de retención.
\end{knowledgebox}

\begin{warningbox}{«Limpiar la escena» puede destruir a la vez la evidencia y la recuperación}
Eliminar cuentas, reinstalar sistemas o rotar todas las credenciales puede ser una containment necesaria, pero si no se registra ni se recopila nada, se pierde la root-cause evidence. A la inversa, retrasar la protección de los usuarios para hacer el análisis forense también puede ampliar el daño. El IC debe hacer que el technical lead y forensics/counsel definan con claridad qué acciones pueden ejecutarse de inmediato, cuáles requieren un snapshot y cuáles no pueden preservarse, pero deben registrarse con su justificación.
\end{warningbox}

\subsection{Decision log: separar hechos, análisis y autorización}

Un registro de calidad no consiste en conservar para siempre cada mensaje de Slack, sino en mantener un decision log estructurado: hechos técnicos actuales, incógnitas, análisis jurídico, opciones de negocio, decision owner, hora, fundamento y condiciones de nueva revisión. El \term{privilege} (privilegio de confidencialidad abogado-cliente), en el contexto de Estados Unidos, suele proteger comunicaciones confidenciales específicas realizadas para obtener o brindar asesoría jurídica; incluir a un abogado en un canal o marcar un documento como privileged no hace que todos los hechos de negocio desaparezcan ni que queden protegidos automáticamente.

\lecturefigure{09-decision-log.png}{Los hechos técnicos, el análisis jurídico, las decisiones de la alta dirección y el registro deben mantener sus límites dentro de una misma cadena.}{Subtítulos locales 21:00--25:40; guías del NIST y la SEC; redibujo conceptual.}

\begin{knowledgebox}{Lectura de la figura: registrar «quién decidió qué», no un consenso ficticio}
Los technical facts deben incluir evidence y confidence; el counsel analysis expone las obligaciones aplicables, los supuestos y el alcance de la confidencialidad; la executive decision enumera las options, la risk acceptance y el approver; record/audit conserva el owner, el rationale, la next review y la superseded decision. Cuando los hechos cambien más adelante, deben añadirse correcciones en lugar de sobrescribir las entradas anteriores. Así se evita que se suponga que el CSO tomó por sí solo una decisión jurídica, y también que los directivos afirmen después que nunca recibieron la información sobre el riesgo.
\end{knowledgebox}

\begin{importantbox}{Plantilla mínima de decision record}
\textbf{Time/owner}; \textbf{known facts and evidence links}; \textbf{unknowns}; \textbf{options}; \textbf{technical recommendation}; \textbf{legal/materiality analysis owner}; \textbf{business decision and approver}; \textbf{external communication}; \textbf{next checkpoint}. No escriba conjeturas en el campo de hechos, ni haga pasar el campo de recommendation por una decisión ya aprobada.
\end{importantbox}

\subsection{Resumen de la sección}

La cadena de evidencia responde «cómo lo sabemos», y el decision log responde «quién tomó qué decisión con base en qué hechos conocidos y desconocidos». Ambos protegen en conjunto la calidad de la investigación, los intereses de los usuarios y la responsabilidad de la organización; de ninguna manera sirven para fabricar memorandos que eximan de responsabilidad a una persona.
\section{Materiality y Disclosure: desde cuándo se cuentan los cuatro días hábiles}

El equipo de seguridad puede evaluar el impacto técnico, pero no puede decidir por sí solo todas las obligaciones de divulgación. La \term{materiality} (importancia relevante), en el contexto de la legislación bursátil de Estados Unidos, se centra en si es sustancialmente probable que un inversionista razonable considere esa información importante para su decisión de inversión; no la determinan por sí solos el número de registros, el CVSS, los minutos de interrupción ni la etiqueta del atacante. Además, cada jurisdicción, sector, contrato y regla de notificación a usuarios tiene su propio trigger, por lo que la organización necesita un parallel analysis en lugar de buscar un «umbral universal».

\subsection{Facts, materiality, audience y timing}

El flujo de trabajo de divulgación establece primero hechos confiables; después, las personas adecuadas analizan el impacto operativo, financiero, reputacional, legal y estratégico, y luego se identifican los destinatarios y los plazos. Los destinatarios pueden incluir a los usuarios afectados, socios de negocio, aseguradoras, autoridades policiales, reguladores de privacidad o del sector, y el mercado público. Las distintas notificaciones pueden tener propósitos y niveles de detalle diferentes, pero no pueden contradecirse. Al leer este flujo conviene notar en particular lo siguiente: la technical severity solo responde con qué rapidez deben asignarse recursos de defensa, mientras que la materiality responde si el incidente es importante para destinatarios legales y de decisión específicos; la primera puede detonar la segunda, pero no puede sustituirla automáticamente.

\lecturefigure{08-disclosure-decision.png}{La decisión de divulgación parte de los hechos técnicos, pasa por el análisis de materiality y llega a la elección de destinatarios, contenido y plazos.}{SEC 2023 cybersecurity disclosure final rule; NIST guidance; rediseño conceptual.}

\begin{knowledgebox}{Lectura de la figura: primero se determina, después se inicia el reloj correspondiente}
Facts enumera scope, data, systems, operational impact y confidence; materiality analiza investor/customer/user harm, financial/operational effect y cumulative impact; audience enumera users, regulators, law enforcement y partners; timing hace después la correspondencia con contract, statute, Form 8-K o approved delay. El orden esencial es «tras el descubrimiento, determinar la importancia relevante sin demora injustificada; una vez determinada como relevante, calcular la ventana de presentación correspondiente», y no que todos los incidentes tengan uniformemente cuatro días a partir de su descubrimiento.
\end{knowledgebox}

\begin{importantbox}{Alcance preciso de la regla de cuatro días hábiles de la SEC}
Para las empresas estadounidenses que cotizan en bolsa y a las que aplica, la regla de la SEC exige determinar la importancia relevante sin demora injustificada después de descubrir el incident; si la empresa determina que el incidente es relevante, normalmente debe presentar el Form 8-K Item 1.05 \textbf{dentro de los cuatro días hábiles posteriores a esa determinación de importancia relevante}. El \term{Form 8-K} es el formulario de reporte corriente con el que las empresas cotizadas informan sobre ciertos eventos importantes. La regla incluye un mecanismo de aplazamiento para cuando el Procurador General de Estados Unidos determine que la divulgación inmediata representaría un riesgo sustancial para la seguridad nacional o la seguridad pública. Su aplicación concreta debe evaluarla un counsel calificado.
\end{importantbox}

\begin{warningbox}{No hay que esperar a que «todos los hechos sean cien por ciento seguros»}
La investigación puede durar semanas, pero el plazo de notificación puede ser más corto. El objetivo de la gobernanza no es usar la incertidumbre para aplazar, sino dejar claro lo conocido, lo desconocido, el nivel de confianza y las actualizaciones previstas. Dar cifras inexactas demasiado pronto daña la confianza, y poner en marcha demasiado tarde el materiality workstream también crea riesgo de incumplimiento. Hay que involucrar al disclosure owner desde las primeras etapas del incidente y establecer checkpoint de reevaluación periódica.
\end{warningbox}

\begin{knowledgebox}{Cómo evita RACI la dilución de responsabilidades}
\term{RACI} es la matriz de responsabilidades Responsible, Accountable, Consulted, Informed. Para el technical scoping, el forensics lead puede ser Responsible y el IC Accountable; para las obligaciones legales, el counsel es Responsible y el GC Accountable; para la divulgación de asuntos operativos importantes, el CEO o el disclosure committee es Accountable. RACI no busca que los cuatro tipos de rol carguen la responsabilidad por partes iguales, sino dejar claro quién es el único responsable final de cada deliverable.
\end{knowledgebox}

\subsection{Resumen de la sección}

Disclosure es una decisión basada en evidencia, interfuncional y sujeta a varios plazos. El equipo de seguridad aporta hechos oportunos y confiables, el counsel interpreta las obligaciones, y la alta dirección o el comité asume las decisiones operativas y de divulgación pública; los «cuatro días» comienzan con la determinación de importancia relevante, no son una cuenta regresiva mecánica a partir de la primera alerta.
\section{Cronología del caso Sullivan: lo que se anticipaba en clase y el estado procesal final}

La clase utiliza el incidente de Uber de 2016 para discutir el bug bounty, la divulgación y la regulation by enforcement. Para mantener la precisión didáctica, es necesario leerlo en cuatro niveles: el relato personal de Sullivan en clase; la teoría del gobierno en la acusación y el juicio; las resoluciones que ya dictaron el jurado y los tribunales; y los controles de ingeniería que este texto extrae de todo ello. Ningún nivel puede sustituir a otro.

\subsection{Hitos procesales de 2016--2026}

El DOJ anunció en 2022 que el jurado emitió un veredicto de culpabilidad por obstrucción de un procedimiento de la FTC y por misprision of felony; en 2023 el tribunal impuso tres años de probation. Cuando se grabó la clase, Sullivan discutió el \term{nexus}: el requisito de vinculación que debe existir, en el contexto de la obstruction, entre la conducta imputada y el procedimiento oficial, y anticipó que el fallo Fischer de la Suprema Corte de 2024 podría favorecer la apelación. Posteriormente, el Noveno Circuito confirmó la condena y publicó una opinión modificada; la Suprema Corte denegó el \term{certiorari} (solicitud de revisión del expediente por el tribunal superior), es decir, se negó a conocer de esa solicitud de apelación, lo cual no significa que la Suprema Corte haya respaldado por separado cada cuestión de fondo.

\lecturefigure{10-current-legal-timeline.png}{La pending appeal mencionada en clase quedó sustituida por los resultados procesales de los tribunales en 2025--2026.}{DOJ; Ninth Circuit No. 23-927; Supreme Court docket 25-1082; redibujo conceptual.}

\begin{knowledgebox}{Lectura de la figura: separar hechos, opiniones y resoluciones}
2016 es el incidente de fondo; la jury conviction de 2022 y la sentencing de 2023 son resultados del juicio; la Ninth Circuit opinion de 2025 es la resolución de apelación; la cert denial de 2026 cierra esta ronda de solicitud de revisión ante la Suprema Corte. Lo que se afirma en clase sobre el nexus, el harmless error o el impacto en la industria corresponde a la postura procesal de enero de 2025. Las notas actuales pueden explicar sus implicaciones de gobernanza, pero no pueden seguir diciendo «apelación pendiente» ni presentar las predicciones del ponente como conclusiones del tribunal.
\end{knowledgebox}

\begin{warningbox}{El caso no es una prueba jurídica general para los VDP}
El caso involucra a una empresa concreta, una investigación de la FTC, los hechos del incidente, comunicaciones y registros. De un solo caso penal no puede deducirse que «pagar un bounty siempre es ilegal» ni que «firmar un NDA siempre es encubrimiento», como tampoco que «basta con tener un VDP para no incurrir en responsabilidad». Lo reutilizable en ingeniería son los mecanismos de clasificación, escalamiento, evidencia, divulgación y toma de decisiones conjunta; las conclusiones jurídicas concretas deben remitirse a la jurisdicción y a los hechos.
\end{warningbox}

\teachervoice{Sullivan veía en clase a Fischer y al nexus como una oportunidad importante para la apelación. La forma correcta de conservar la teacher voice es registrar «por qué lo juzgaba así en ese momento» y, al mismo tiempo, dejar claro que los resultados judiciales posteriores sustituyeron el estado pending; la fidelidad no consiste en congelar un momento histórico como si fuera un hecho actual.}

\begin{importantbox}{Cuatro etiquetas para las afirmaciones}
Al redactar un case study, distinga de forma explícita: \textbf{speaker account} (lo que afirma el ponente), \textbf{government allegation} (lo que sostiene el gobierno), \textbf{adjudicated fact/outcome} (lo que resolvieron el jurado o los tribunales) y \textbf{engineering judgment} (lo que recomienda este texto). Esto reduce las conjeturas sobre las motivaciones de las personas y permite al lector saber qué parte puede adoptarse directamente en el diseño de sistemas.
\end{importantbox}

\subsection{Resumen de la sección}

El valor didáctico de esta cronología está en corregir la versión anterior: al 11 de agosto de 2026, la apelación ya no está pending. Más importante aún, una organización no puede tratar el desenlace jurídico como una variable externa que solo importa después del incidente; el registro de los hechos, la asignación de responsabilidades y las vías de divulgación deben diseñarse antes de que ocurra el incidente.
\section{Executive Accountability: responsabilidad compartida, pero no diluida}

Los incidentes de seguridad graves suelen dejar al descubierto la estructura de la organización: quién recibe los hechos, quién puede exigir más investigación, quién aprueba los pagos, quién decide no notificar, quién informa al consejo de administración. Cargar toda la responsabilidad sobre el CSO fomenta una cultura defensiva. Repartirla por igual entre una decena de comités, en cambio, hace que nadie sea realmente accountable. Una gobernanza madura necesita que coexistan la shared responsibility y un named decision owner.

\subsection{Los decision rights del CSO, el GC, el CEO y el board}

El CSO debe asegurarse de que el riesgo se escale con precisión. El GC debe asegurarse de que un equipo calificado atienda las cuestiones legales. El CEO debe responder por los recursos y por las decisiones de negocio importantes. El board o el disclosure committee debe aportar supervisión y challenge. El CSO no puede sustituir con un «legal approved» su propia obligación de veracidad técnica. Tampoco el GC puede juzgar el impacto por su cuenta cuando falta evidence. Y menos aún debe el CEO oír hablar por primera vez de la arquitectura de seguridad justo antes de la divulgación pública.

\lecturefigure{12-executive-accountability.png}{La responsabilidad en ciberseguridad la comparten el CSO, el GC, el CEO y el board, pero cada decisión necesita un owner explícito.}{Subtítulos locales 25:20--35:20; conceptos de gobernanza redibujados.}

\begin{knowledgebox}{Lectura de la figura: se comparte la información, no una responsabilidad difusa}
El CSO aporta el technical scope, la security recommendation y los escalation facts. El GC aporta los legal duties, la regulator strategy y el privilege advice. El CEO toma las decisiones de enterprise risk y de disclosure. El board o el committee supervisa, cuestiona y exige una revisión posterior. Los límites de la figura no son silos: cada rol puede hacer challenge a los demás workstreams, pero no puede trasladar a otro su propio deliverable. El registro final debe poder responder quién es responsable de los hechos, quién del análisis y quién aprobó las acciones.
\end{knowledgebox}

\teachervoice{El ponente recomienda que los líderes de seguridad construyan la relación con el CEO y el GC antes de una crisis, en lugar de pedirles que entiendan un incident complejo en la primera llamada importante. Nota de clase: la executive readiness es una interface que se entrena de antemano y requiere brief periódicos, tabletop y un umbral explícito de «cuándo despertarme».}

\begin{warningbox}{«Ya se lo había dicho» no significa que la gobernanza esté completa}
Un solo correo o un aviso verbal no demuestra que el destinatario entienda el riesgo, tenga el decision right o haya completado las acciones siguientes. Un escalamiento de alto riesgo debe incluir severity, evidence, unknowns, recommended actions, deadline, un ask explícito y una confirmación de recepción. Si la decisión se aplaza o se rechaza, deben registrarse el owner, el motivo y la siguiente revisión.
\end{warningbox}

\begin{importantbox}{La forma correcta de proteger a los líderes de seguridad}
Proteger no significa redactar memos para cubrirse. La organización debe proporcionar una job authority clara, un escalamiento por escrito, asesoría legal independiente, acceso al consejo de administración, D\&O/indemnification y el employment support adecuado. Por su parte, el líder debe mantenerse honesto respecto de los hechos, no excederse en sus atribuciones emitiendo conclusiones legales, negarse a eliminar registros desfavorables y, ante desacuerdos importantes, recurrir a los canales formales de escalamiento.
\end{importantbox}

\subsection{Resumen de la sección}

El objetivo de la accountability es que decidan las personas que tienen la capacidad, la información y la autoridad para hacerlo. El CSO no debe convertirse en el único responsable de todos los riesgos de seguridad, pero tampoco puede renunciar, con el pretexto de una «decisión colectiva», a escalar con precisión ni a su juicio profesional.
\section{Resilience y Public Service: cómo seguir trabajando después de un incidente}

Los incidentes de seguridad y los procesos legales afectan la salud, la identidad y la carrera de empleados, investigadores, usuarios y líderes. La \term{resilience} (resiliencia) no consiste en fingir que no hubo daño, sino en reconstruir de forma gradual una capacidad de servicio sostenible mediante apoyo, recuperación, reflexión y reincorporación. El curso menciona el trabajo de Sullivan en un proyecto de servicio público en Ucrania; esta parte no puede tomarse como prueba del caso, pero aporta una perspectiva importante sobre el sistema humano.

\subsection{Human recovery también es parte del incident lifecycle}

El postmortem tradicional se centra en la root cause, el control gap y el owner, pero puede pasar por alto las guardias prolongadas, la presión pública, el daño moral y la carga familiar. Si una organización solo exige «restablecer el servicio», agota al personal clave antes del siguiente incidente. El recovery plan debe incluir rotación, apoyo psicológico, asistencia legal y de HR, licencias, ajustes de funciones y una vía para reincorporarse al trabajo.

\lecturefigure{14-resilience-loop.png}{La resiliencia surge de un ciclo de apoyo, recuperación, reconstrucción del sentido y regreso al servicio.}{Subtítulos locales 27:00--29:20; redibujo conceptual.}

\begin{knowledgebox}{Lectura de la figura: recuperarse no es volver al estado previo al incidente}
El incident/case provoca un impacto en la operación y en las personas; el support aporta family, peer, professional y legal care; la recovery requiere rest, perspective y new routines; service/learning convierte la experiencia en community work, mentoring y better systems. Las flechas no garantizan que cada persona se recupere de forma lineal, sino que recuerdan a la organización que debe dejar margen para ritmos distintos. Interpretar la resilience como «la persona debe ser más fuerte» oculta la responsabilidad institucional.
\end{knowledgebox}

\begin{warningbox}{No usar el servicio público ni el crecimiento personal como blanqueo legal}
Las buenas acciones de una persona después de un incidente pueden explicar sus decisiones de valores y su proceso de recuperación, pero no modifican por sí mismas los hechos ya juzgados; del mismo modo, la responsabilidad legal no debe llevar a la organización a descuidar la salud y la dimensión humana. Las notas del curso deben permitir que ambos juicios coexistan, en lugar de sustituir el análisis de la evidencia por un relato inspirador.
\end{warningbox}

\subsection{Government technical talent}

El sector público necesita talento técnico que entienda cloud, identity, AI, supply chain e incident operations; de lo contrario, las normas pueden describir objetivos sin poder evaluar su implementación. Que los líderes técnicos participen en el government no significa solo ocupar un cargo de tiempo completo: también incluye standards work, advisory groups, comentarios públicos, incident coordination y periodos cortos de servicio. Los participantes deben gestionar los conflictos de interés y anteponer la responsabilidad pública al lobbying de su empresa.

\lecturefigure{15-government-talent.png}{El talento técnico solo mejora la calidad de la gobernanza pública si permanece dentro del ciclo de retroalimentación entre políticas, implementación e industria.}{Subtítulos locales 36:20--39:17; redibujo conceptual.}

\begin{knowledgebox}{Lectura de la figura: la circulación del talento debe generar memoria institucional}
La technical expertise ayuda a describir systems, threats y operational cost; el policy design traduce los objetivos en rules, incentives y oversight; la implementation construye agency capability, metrics y enforcement; el industry feedback devuelve evidence, failure modes y burden. Si los expertos solo «aterrizan» por poco tiempo sin documentación, career path ni permanent staff, el conocimiento se va con cada persona; si la industria solo aporta posturas y no datos verificables, tampoco pueden formarse normas de alta calidad.
\end{knowledgebox}

\teachervoice{Al final, el ponente pidió a los profesionales técnicos que no consideren al government como un problema ajeno. Nota de clase: cuando los sistemas de software sostienen la vida pública, quienes entienden esos sistemas deben ayudar a diseñar las restricciones públicas; de lo contrario, las normas aparecerán de todos modos, solo que con mayor probabilidad en condiciones de información insuficiente y bajo la presión de una crisis.}

\subsection{Resumen de la sección}

La resilience incorpora la recuperación personal a la capacidad de seguridad a largo plazo, y el public service devuelve la experiencia técnica a la retroalimentación institucional. Ambos exigen pasar del relato heroico a los sistemas de apoyo y a la memoria organizacional.
\section{Ejercicio integrador: del VDP Report a la divulgación y el aprendizaje}

Para cerrar, un tabletop enlaza los mecanismos de esta clase en una ruta que puede ensayarse. Supongamos que un investigador envía por el VDP un reporte de alto impacto; a primera vista parece solo una vulnerabilidad, pero durante la verificación se descubre el uso real de una credential y un posible acceso a datos. El objetivo del ejercicio no es adivinar los motivos del atacante, sino observar si el equipo puede escalar a tiempo, preservar la evidencia, proteger a los usuarios, completar el análisis de materiality y mantener una comunicación coherente.

\subsection{Tabletop en cinco etapas}

En cada etapa solo se entrega a los participantes una parte de los hechos, y se les pide registrar sus supuestos, el owner y el siguiente paso. El Facilitator debe introducir incertidumbre a propósito: registros faltantes, dependencias de terceros, un investigador que exige respuesta rápida, un área de negocio que quiere mantener el servicio y medios de comunicación que empiezan a hacer preguntas. Una buena respuesta no es «cerrar todo de inmediato» ni «esperar a que termine la investigación», sino una que explique el beneficio, el costo, la evidencia requerida y quién aprueba cada acción.

\lecturefigure{16-tabletop.png}{El tabletop enlaza VDP, incident, materiality, notification y learning en un ejercicio de extremo a extremo.}{Diseño integrador de esta clase; concepto redibujado.}

\begin{knowledgebox}{Lectura de la figura: cada etapa debe entregar un resultado con capacidad de auditoría}
El VDP report entrega el case ID, el scope y el acknowledgement; el incident entrega la severity, el IC, la containment y el evidence plan; la governance entrega el materiality owner, la RACI y el decision log; notify/learn entrega mensajes coherentes, el deadline, la recovery y los remediation owners. El éxito del ejercicio no se mide por «no haber discutido», sino por si la discusión giró en torno a los mismos hechos, si se escaló antes del deadline y si la decisión la tomó el rol correcto.
\end{knowledgebox}

\begin{importantbox}{Tabletop injects}
\begin{enumerate}
  \item El reporte proviene de un out-of-country researcher, describe una vulnerabilidad in-scope, pero adjunta registros reales de usuarios;
  \item los internal logs muestran que la credential ya se había usado, en una fecha anterior al reporte;
  \item un SaaS de terceros también podría estar afectado, con una ventana contractual de notificación distinta;
  \item la corrección requiere una breve interrupción del servicio, y el responsable de negocio pide posponerla;
  \item el counsel solicita un análisis de materialidad, y communications recibe preguntas de un periodista;
  \item nueva evidencia contradice la cifra inicial de registros, y el equipo debe corregir, no sobrescribir, la decisión anterior.
\end{enumerate}
\end{importantbox}

\begin{warningbox}{El ejercicio no debe buscar una «respuesta correcta única»}
Las obligaciones reales dependen de los hechos y de la jurisdicción. El tabletop debe poner a prueba las interfaces: si se identifica el trigger de escalamiento, si se preserva la evidence clave antes y después de la containment, si el GC, el CEO o el committee se incorporan cuando hace falta, y si se pueden explicar las decisiones ante los usuarios y el público. Si el ejercicio solo evalúa la memorización de políticas, no podrá revelar los puntos de ruptura de la organización.
\end{warningbox}

\subsection{Resumen de la sección}

El ejercicio de extremo a extremo convierte las políticas, de documentos, en músculo organizacional. Solo cuando VDP, incident command, evidence, materiality, executive decision y learning colaboran en un mismo escenario puede el equipo descubrir los vacíos de responsabilidad que existen en la realidad.
\section{Síntesis y ampliación}

Esta clase replantea el security leadership: deja de ser la pregunta «¿es el CSO lo bastante valiente?» y se convierte en un diseño organizacional verificable:

\begin{enumerate}
  \item \textbf{La seguridad es un operating system}: people, process, platform y governance determinan juntos el resultado;
  \item \textbf{Las reglas públicas necesitan retroalimentación}: las empresas deben llevar la evidence de sistemas reales a la elaboración de políticas y, al mismo tiempo, aceptar la accountability pública;
  \item \textbf{El VDP gestiona el canal de reportes}: scope, acknowledgement, corrección y divulgación coordinada conforman un contrato de confianza;
  \item \textbf{El Bounty gestiona incentivos opcionales}: una recompensa, un NDA o un payment no pueden autorizar a posteriori una conducta dañina;
  \item \textbf{El Incident command gestiona la ejecución interfuncional}: el IC sincroniza los flujos de trabajo, pero no monopoliza los juicios técnicos, legales ni de negocio;
  \item \textbf{La Evidence y el decision log garantizan que todo pueda revisarse}: hechos, análisis, decisiones y correcciones deben registrarse en capas separadas;
  \item \textbf{Materiality no equivale a severity}: para las empresas cotizadas a las que aplica, el plazo de cuatro días hábiles del Form 8-K suele contarse a partir de que se determina la importancia del incidente;
  \item \textbf{El estado de los casos debe actualizarse}: la pending appeal mencionada en clase quedó superada por la sentencia del Noveno Circuito de 2025 y por la cert denial de 2026;
  \item \textbf{La Accountability se comparte, pero no se diluye}: CSO, GC, CEO y board tienen cada uno un deliverable y un decision right claramente definidos;
  \item \textbf{La Resilience y el public service son capacidades de largo plazo}: recuperar al personal, conservar la memoria institucional y mejorar la gobernanza pública de la tecnología.
\end{enumerate}

\begin{importantbox}{Tarea de Security Governance}
Diseña un playbook de extremo a extremo para una empresa SaaS ficticia: publica el VDP y los límites del safe-harbor; define los VDP-to-incident triggers; traza la matriz RACI de IC, technical, legal/comms y business; proporciona plantillas de evidence chain y de decision-log; enumera los workstreams de divulgación a customer, regulator, law-enforcement y public-market; por último, diseña seis tabletop injects. Cada decisión clave debe indicar owner, evidence, deadline y review condition.
\end{importantbox}

\begin{knowledgebox}{Autoevaluación final: cinco pares que no deben confundirse}
No confundas vulnerability con incident, ni severity con legal materiality, ni payment con authorization, ni speaker account con adjudicated outcome, ni shared responsibility con diffuse accountability. Mientras estos cinco límites estén claros, es más probable que la organización mantenga íntegros los hechos y las responsabilidades bajo presión.
\end{knowledgebox}

\subsection{Lecturas adicionales}

Los materiales siguientes no están en el mismo nivel: CISA, DOJ y NIST ayudan principalmente a diseñar el VDP y el incident-response process; los documentos de la SEC explican un marco específico de divulgación en mercados públicos; las opiniones judiciales y el docket registran el resultado procesal de los casos. Al leerlos, primero determina si quieres responder «cómo construir controles», «qué obligaciones aplican» o «qué juicio emitió el tribunal sobre el expediente existente»; así evitarás sustituir el análisis jurídico por guías de ingeniería, y también sustituir principios generales de ingeniería por un solo caso.

{\small
\begin{itemize}[nosep,leftmargin=*]
  \item \href{https://www.cisa.gov/vulnerability-disclosure-policy-template}{\textbf{CISA Vulnerability Disclosure Policy Template}}: ejemplos de scope, reporting, expectations y coordinated disclosure.
  \item \href{https://www.cisa.gov/news-events/directives/bod-20-01-develop-and-publish-vulnerability-disclosure-policy}{\textbf{CISA BOD 20-01}}: contexto de política del VDP para las agencias federales.
  \item \href{https://www.justice.gov/criminal/criminal-ccips/page/file/983996/dl}{\textbf{DOJ Framework for a Vulnerability Disclosure Program}}: marco sobre autorización, good faith y consideraciones de aplicación de la ley.
  \item \href{https://csrc.nist.gov/pubs/sp/800/61/r3/final}{\textbf{NIST SP 800-61 Rev. 3}}: integra la incident response en las actividades de gestión de riesgos del CSF 2.0.
  \item \href{https://www.sec.gov/files/33-11216-fact-sheet.pdf}{\textbf{SEC Cybersecurity Disclosure Final Rule Fact Sheet}}: panorama del material incident y de los plazos del Form 8-K.
  \item \href{https://cdn.ca9.uscourts.gov/datastore/opinions/2025/11/12/23-927.pdf}{\textbf{United States v. Sullivan, No. 23-927}}: opinión modificada del Noveno Circuito.
  \item \href{https://www.supremecourt.gov/docket/docketfiles/html/public/25-1082.html}{\textbf{Supreme Court Docket 25-1082}}: registro de la certiorari petition y de su rechazo del 29 de junio de 2026.
\end{itemize}
}

\begin{knowledgebox}{Orden de lectura sugerido: de los controles ejecutables a los límites judiciales}
Primero usa el CISA template para revisar si la policy pública establece scope, compromisos y canales de contacto; después usa el DOJ framework para examinar los límites de good-faith y de autorización; luego usa NIST SP 800-61 Rev. 3 para escalar el reporte a una incident response de nivel organizacional; si la organización es una empresa cotizada estadounidense a la que aplica la regla, lee el SEC fact sheet y el texto original de la final rule; por último, contrasta la Ninth Circuit opinion con el Supreme Court docket para confirmar el estado procesal del caso visto en clase. En cada paso registra la fecha de la fuente y su ámbito de aplicación, porque las políticas, las normas y el estado de los casos seguirán cambiando.
\end{knowledgebox}

\begin{importantbox}{Ejercicio de Source triangulation}
Elige un incident hipotético y escribe por separado: las recomendaciones operativas que el equipo técnico obtiene de NIST, las obligaciones de notificación y divulgación que counsel debe verificar, las decisiones de negocio que debe tomar la dirección y las preguntas que los materiales judiciales \textbf{no pueden} responder por ti. Si la misma conclusión aparece en las cuatro columnas, eso suele indicar que los niveles de evidencia siguen mezclados.
\end{importantbox}

\end{document}
